\documentclass[11pt]{report}
\usepackage[T1]{fontenc}
\usepackage{multirow}
\usepackage{longtable}
\usepackage{arydshln}
\usepackage{amsmath}
\usepackage{amsthm}
\usepackage{amssymb}
\usepackage{hyperref}
\usepackage{xcolor}
\usepackage{mathtools}
\usepackage{sourcecodepro}

\setcounter{secnumdepth}{3}
\setcounter{tocdepth}{3}

\renewcommand{\arraystretch}{1.5} % increase space between table rows

\newcommand{\set}[1]{\left\{#1\right\}}
\newcommand{\infrule}[2]{\frac{\begin{gathered}#1\end{gathered}}{\begin{gathered}#2\end{gathered}}}
\newcommand{\oprule}[2]{\frac{\begin{gathered}#1\end{gathered}}{\begin{gathered}#2\end{gathered}}}
% \newcommand{\infrule}[2]{\colorbox{gray!2}{$\frac{\begin{gathered}#1\end{gathered}}{\begin{gathered}#2\end{gathered}}$}}
\newcommand{\tuple}[1]{\left<#1\right>}
\newcommand{\code}[1]{\texttt{\colorbox{gray!10}{#1}}}
\newcommand{\tuplecode}[1]{\code{<#1>}}
\newcommand{\narr}{<:}
\newcommand{\env}[1]{[\![#1]\!]}
\newcommand{\envConst}{\Pi}
\newcommand{\envVar}{\Upsilon}
\newcommand{\envStructs}{\Sigma}
\newcommand{\envFunc}{\Phi}
\newcommand{\envLoc}{\Lambda}
\newcommand{\envStore}{\Psi}

\theoremstyle{definition}
\newtheorem{definition}{Definition}[section]

\newcommand{\comment}[1]{\parbox{0.9\textwidth}{\color{black!30}\small\sffamily #1}}

\makeatletter
\providecommand{\bigsqcap}{%
  \mathop{%
    \mathpalette\@updown\bigsqcup
  }%
}
\newcommand*{\@updown}[2]{%
  \rotatebox[origin=c]{180}{$\m@th#1#2$}%
}
\makeatother

\begin{document}

\setlength{\doublerulesep}{0.40cm}
% \setlength{\baselineskip}{11pt}
\abovedisplayshortskip=8pt
\belowdisplayshortskip=8pt

\tableofcontents

\newpage

\chapter{Types}

\section{Conversion}

\resizebox{0.5\linewidth}{!}{
	\begin{tabular}{ll}
		Symbol & Meaning \\\hline
		$=$ & Identity \\\hline
		$w$ & Wideneable \\\hline
		$n$ & Narrowable \\\hline
		$W$ & Wideneable if components are wideneable \\\hline
		$N$ & Narrowable if components are narrowable \\\hline
		$!$ & Explicit cast is required (it may fail) \\\hline
	\end{tabular}
}

\subsection{Primitive types}

\resizebox{0.9\linewidth}{!}{
	\begin{tabular}{|c|c|c|c|c|c|c|c|}
		\hline
		$_{from} \quad ^{to}$ & \texttt{Void} & \texttt{Type} & \texttt{Bool} & \texttt{Natural} & \texttt{Int} & \texttt{Real} & \texttt{Text} \\\hline
		\texttt{Void}           & $=$         & $w$         & $w$         & $w$            & $w$        & $w$         & $w$         \\\hline
		\texttt{Type}           &             & $=$         &             &                &            &             &             \\\hline
		\texttt{Bool}           &             &             & $=$         &                &            &             &             \\\hline
		\texttt{Natural}        &             &             &             & $=$            & $w$        & $w$         &             \\\hline
		\texttt{Int}            &             &             &             & $n!$           & $=$        & $w$         &             \\\hline
		\texttt{Real}           &             &             &             & $n!$           & $n!$       & $=$         &             \\\hline
		\texttt{Text}           &             &             &             &                &            &             & $=$         \\\hline
	\end{tabular}
}


\subsection{Non primitive types}

\resizebox{0.9\linewidth}{!}{
	\begin{tabular}{|c|c|c|c|c|c|c|c|}
		\hline
		$_{from} \quad ^{to}$ & \texttt{Void} & \texttt{List} & \texttt{Map} & \texttt{Function} & \texttt{Tuple} & \texttt{Record} & \texttt{Struct} \\\hline
		\texttt{Void}           & $=$         & $w$         & $w$         & $w$            & $w$          & $w$           & $w$           \\\hline
		\texttt{List}           &             & $N! W$      &             &                &              &               &               \\\hline
		\texttt{Map}            &             &             & $N! W$      &                &              &               &               \\\hline
		\texttt{Function}       &             &             &             & $N! W$         &              &               &               \\\hline
		\texttt{Tuple}          &             &             &             &                & $N! W$       & $W$           & $W$          \\\hline
		\texttt{Record}         &             &             &             &                & $N!$         & $N! W$        & $W$           \\\hline
		\texttt{Struct}         &             &             &             &                & $N!$         & $N!$          & $N! W!$       \\\hline
	\end{tabular}
}

\chapter{Inference rules}

\section{Notations and Terminology}

\subparagraph{Judgements and terms} To describe the type system, its types and
the inference rules, this document uses "judgements" on "terms".\\*
For example, $t:T$ means $t$ is a term of type $T$.

\subparagraph{Assumptions} Zero or more assumptions (comma separated) can be
used for a judgement; they're specified before the "turnstile" symbol
$\vdash$.\\*
For example, $\Gamma \vdash t:T$ can be read as "if $\Gamma$ is true, then $t$
is a term of type $T$".

\subparagraph{Premise} Judgements can be made based on zero or more judgements
(premise, space separated).\\*
For example, $\infrule{\Gamma\vdash t_1:T \quad \Delta\vdash
t_2:T}{\Gamma,\Delta\vdash \code{$t_1$ + $t_2$}: T}$ means that with a premise
that $\Gamma\vdash t_1:T$ and $\Delta\vdash t_2:T$ are true, then
$\Gamma,\Delta\vdash \code{$t_1$ + $t_2$}: T$ is true.

\subparagraph{Code} Text written with \code{monospaced} font denotes the text
as written in code verbatim.\\*
Some types may also be written with \texttt{Monospace} font (usually in pascal
case)

\subparagraph{Functions} Functions are denoted by a list of zero or more input
arguments followed by an output argument.\\*
For example, $(\tau_1, \tau_2) \to \rho$ describes a function that takes two
arguments of type $\tau_1$ and $\tau_2$ respectively, and produces a value of
type $\rho$

\subparagraph{"Generic" types} Some types are "generic" (a.k.a. existential
types), meaning that they declare one or more type variable (a.k.a. type
parameters).\\*
The notation $T[\tau:=\sigma]$ describes a generic type $T$ that
takes a type parameter $\tau$ of some type $\sigma$.\\*
The notation $t : \exists\sigma.T$ means that the $t$ is a term of concrete
type $T$ with the type parameter $\tau$ equal to $\sigma$

\comment{TODO: tipi esistenziali? Rivedere notazione}

\subparagraph{Identifiers set} A set $\mathcal{I}=\set{i\mid
i:=\text{identifier}}$ is used to encompass all possible identifiers allowed by
the grammar.

\comment{TODO: formalizzare}

\paragraph{Product types}

\subparagraph{Tuple} A tuple is denoted with the cartesian product operator.\\*
For example, $\tau \times \sigma \times \rho$ describes a product 3-tuple with
three members of type $\tau, \sigma, \rho$.

\subparagraph{Record} The notation $\tuple{a:\tau, b:\sigma}$ describes a
record (a tuple whose members are named)

\subparagraph{Struct} The notation $(\iota, \tuple{a:\tau, b:\sigma})$
describes a struct (a pair of an identifier $\iota$ and a record)

\comment{TODO: rivedere notazione struct}

\begin{definition}[Types widening/narrowing]
	A type $N$ is said to be "narrower" than a type $W$ iff a value described
	by $N$ can also be described by $W$ without any loss of information.\\*
	Conversely, a type $W$ is said to be "wider" than a type $N$ iff it can
	describe any value described by $N$ without loss of information.\\*
	A "narrower than" relationship is denoted using the $\narr$ operator 
\end{definition}

The $\sqcap$ operator produces the least wide type for the operand types

\[
	\infrule{
		\Gamma\vdash \tau_1:\texttt{Type} \quad \tau_2:\texttt{Type} \hfill \\*
		\Gamma\vdash \exists_{\omega:\texttt{Type}}{\tau_1\narr\omega \land \tau_2\narr\omega} \hfill \\*
		\Gamma\vdash \nexists_{\sigma:\texttt{Type}}{\tau_1\narr\sigma \land \tau_2\narr\sigma \land \sigma\narr\omega}
	}{
		\Gamma\vdash \tau_1 \sqcap \tau_2 = \omega
	}
\]

\subsection{Type environment}

The definition of an expression type is done within a "context" (a.k.a. type
environment) which describes the type of any identifier or any other
expressions appearing in the described expression.

The general type environment is denoted by $\Gamma$.

To avoid ambiguities and to more clearly describe expressions, $\Gamma$ is
partioned in the $\envConst,\envVar,\envStructs,\envFunc$ environments:

\begin{itemize}
	\item $\envConst$: the Constants environment
		\[
			\begin{split}
				\envConst: \mathcal{I} & \to \texttt{Type}\\*
				\envConst(i) & \mapsto \tau
			\end{split}
		\]
		A function that maps a constant identifier $i$ to its type $\tau$
	\item $\envVar$: the Variable environment
		\[
			\begin{split}
				\envVar: \mathcal{I} & \to \texttt{Type}\\*
				\envVar(i) & \mapsto \tau
			\end{split}
		\]
		A function that maps a variable identifier $i$ to its type $\tau$
	\item $\envStructs$: the Structs environment
		\[
			\begin{split}
				\envStructs: \mathcal{I} & \to (\mathcal{I}, \texttt{Record})\\*
				\envStructs(i) & \mapsto (i, \rho)
			\end{split}
		\]
		A function that maps a typename $i$ (struct definition identifier) to
		its type definition (identifier $i$ and record $\rho$)
	\item $\envFunc$: the Functions environment
		\[
			\begin{split}
				\envFunc: \mathcal{I} & \to (\mathcal{I}, \texttt{Function})\\*
				\envFunc(i) & \mapsto (i, \phi)
			\end{split}
		\]
		A function that maps a function identifier $i$ to its definition
		signature $(i, \phi)$\\*
		With a type $\tau$, it maps to the function defined within the type
		$\tau$
		\[
			\begin{split}
				\envFunc: (\texttt{Type},\mathcal{I}) & \to (\mathcal{I}, \texttt{Function})\\*
				\envFunc(\tau, i) & \mapsto (i, \phi)
			\end{split}
		\]
\end{itemize}

\comment{TODO: formalizzare structs e funzioni}

\paragraph{Environment extension} Given how scopes and symbols
lookup work, it will be useful to define "new" type
environments extending an existing environment. The notation
$E\env{t:T}$ will be used to denote that the environment $E$ is
extended with $\env{t:T}$.

Using the $\envVar$ variables environment as example

\[
	\envVar\env{t:T}(u) \mapsto \begin{cases}
		T &\text{ if } u = t\\
		\envVar(u) &\text{ if } u \neq t\\
	\end{cases}
\]

\newpage

\section{Types}

\comment{TODO: menzionare uso di type system strutturale e nominale}\\*
\comment{TODO: usare 0-based indexes}\\*
\comment{TODO: usare notazione tipi polimorfici invece che solo Type}

\subsection{Primitive types}

\[
	\begin{split}
		&\Gamma\vdash \texttt{Void}:\texttt{Type} \\
		&\Gamma\vdash \texttt{Bool}:\texttt{Type} \\
		&\Gamma\vdash \texttt{Natural}:\texttt{Type} \\
		&\Gamma\vdash \texttt{Int}:\texttt{Type} \\
		&\Gamma\vdash \texttt{Real}:\texttt{Type} \\
		&\Gamma\vdash \texttt{Text}:\texttt{Type} \\
	\end{split}
\]

\subsection{Composite types}

\[
	\begin{split}
		&\Gamma\vdash \texttt{List}[\sigma:=\texttt{Type}]:\texttt{Type} \\
		&\Gamma\vdash \texttt{Map}[\kappa:=\texttt{Type},\nu:=\texttt{Type}]:\texttt{Type} \\
	\end{split}
\]

\[
	\infrule{
		\Gamma\vdash \tau:\texttt{Type}
	}{
		\Gamma\vdash \code{[$\tau$]}: \texttt{List}[\sigma := \tau]
	}
\]

\[
	\infrule{
		\Gamma\vdash \tau:\texttt{Type} \quad \sigma:\texttt{Type}
	}{
		\Gamma\vdash \code{[$\tau$:$\sigma$]}: \texttt{Map}[\kappa:=\tau, \nu:=\sigma]
	}
\]

\subsection{Product types}

\[
	\text{Tuple}\qquad
	\infrule{
		\Gamma\vdash \forall_{i\in[1,n]}{\tau_i:\texttt{Type}} \quad n \geq 0
	}{
		\Gamma\vdash \tuplecode{$\tau_1,...,\tau_n$} : \tau_1 \times ... \times \tau_n
	}
\]

\[
	\text{Record}\qquad
	\infrule{
		\Gamma\vdash n \geq 0 \quad \forall_{i\in[1,n]}{\tau_i:\texttt{Type}, a_i \in \mathcal{I}}
	}{
		\Gamma\vdash \tuplecode{$a_1\tau_1,...,a_n\tau_n$} : \tuple{a_1:\tau_1, ..., a_n:\tau_n}
	}
\]

\[
	\text{Record (empty)} \qquad
	\Gamma\vdash \tuplecode{:} : \tuple{}
\]

\[
	\text{Struct (user defined)} \qquad
	\infrule{
		\Gamma\vdash \eta \in \mathcal{I} \quad \rho:\texttt{Record}
	}{
		\Gamma\vdash \code{$\eta$ $\rho$}: (\eta,\rho)
	}
\]

\comment{TODO: che notazione usare per \texttt{Struct}???}

\subsection{Functions}

\[
	\infrule{
		\Gamma\vdash n \in \mathbb{N} \quad \forall_{i\in[1,n]}{\tau_i:\texttt{Type},a_i\in\mathcal{I}} \quad \rho:\texttt{Type}
	}{
		\Gamma\vdash \code{::($a_1 \tau_1,...,a_n \tau_n$) => $\rho$} : (\tau_1,...,\tau_n) \to \rho
	}
\]

\[
	\infrule{
		\Gamma\vdash n \in \mathbb{N} \quad \forall_{i\in[1,n]}{\tau_i:\texttt{Type},a_i\in\mathcal{I}}
	}{
		\Gamma\vdash \code{::($a_1 \tau_1,...,a_n \tau_n$)} : (\tau_1,...,\tau_n) \to \texttt{Void}
	}
\]

\subsection{Range}

\comment{TODO: tipo concreto non implementato ancora}

\subsection{Compatibility rules}

\[
	\infrule{
		\Gamma\vdash \tau : \texttt{Type}
	}{
		\Gamma\vdash \texttt{Void} \narr \tau
	}
\]

\[
	\infrule{
		\Gamma\vdash \tau:\texttt{Type}
	}{
		\Gamma\vdash \tau \narr \tau
	}
	\qquad \text{Reflexivity}
\]

\[
	\infrule{
		\Gamma\vdash \tau_1 \narr \tau_2 \quad \tau_2 \narr \tau_3
	}{
		\Gamma\vdash \tau_1 \narr \tau_3
	}
	\qquad \text{Transitivity}
\]

\[
	\Gamma\vdash \texttt{Natural} \narr \texttt{Int}
	\qquad \Gamma\vdash \texttt{Int} \narr \texttt{Real}
\]

\paragraph{Composite types}

\[
	\infrule{
		\Gamma\vdash \tau_1:\texttt{List}[\sigma_1:=\texttt{Type}]\\*
		\Gamma\vdash \tau_2:\texttt{List}[\sigma_2:=\texttt{Type}]\\*
		\Gamma\vdash \sigma_1 \narr \sigma_2
	}{
		\Gamma\vdash \tau_1 \narr \tau_2
	}
\]

\[
	\infrule{
		\Gamma\vdash \tau_1:\texttt{Map}[\kappa_1:=\texttt{Type}, \nu_1=\texttt{Type}]\\*
		\Gamma\vdash \tau_2:\texttt{Map}[\kappa_2:=\texttt{Type}, \nu_2=\texttt{Type}]\\*
		\Gamma\vdash \kappa_1 \narr \kappa_2 \land \nu_1 \narr \nu_2
	}{
		\Gamma\vdash \tau_1 \narr \tau_2
	}
\]

\newpage

\paragraph{Product types}

\subparagraph{Tuple}

\[
	\infrule{
		\Gamma\vdash n \geq 0 \quad \forall_{i\in[1,n]}{\tau_i:\texttt{Type}} \quad \theta_1:\bigtimes_{i \in [1,n]}{\tau_i}\\*
		\Gamma\vdash m \geq n \quad \forall_{i\in[1,m]}{\omega_i:\texttt{Type}} \quad \theta_2:\bigtimes_{j \in [1,m]}{\omega_j}\\*
		\Gamma\vdash \forall_{i \in [1,n]}{\tau_i \narr \omega_i}
	}{
		\Gamma\vdash \theta_1 \narr \theta_2
	}
\]

A tuple type $\theta_1$ is narrower than the tuple type $\theta_2$ if the first $n$
fields of $\theta_1$ have a type $\tau$ narrower than the first $n$ fields
types $\omega$ of $\theta_2$

\vspace{\baselineskip}

\[
	\infrule{
		\Gamma\vdash n \geq 0 \quad \forall_{i\in[1,n]}{\tau_i:\texttt{Type}} \quad \theta:\bigtimes_{i \in [1,n]}{\tau_i}\\*
		\Gamma\vdash m \geq n \quad \forall_{j\in[1,m]}{\eta_j \in \mathcal{I}, \omega_j:\texttt{Type}} \quad \rho:\tuple{\eta_1:\omega_1,...,\eta_m:\omega_m} \\*
		\Gamma\vdash \forall_{i \in [1,n]}{\tau_i \narr \omega_i}
	}{
		\Gamma\vdash \theta \narr \rho
	}
\]

A tuple type $\theta$ is narrower than the record type $\rho$ if the first $n$
fields of $\theta$ have a type $\tau$ narrower than the first $n$ fields type
$\omega$ of $\rho$ (fields names are ignored)

\vspace{\baselineskip}

\[
	\infrule{
		\Gamma\vdash n \geq 0 \quad \forall_{i\in[1,n]}{\tau_i:\texttt{Type}} \quad \theta:\bigtimes_{i \in [1,n]}{\tau_i}\\*
		\Gamma\vdash \iota \in \mathcal{I} \quad m \geq n \quad \forall_{j\in[1,m]}{\eta_j \in \mathcal{I}, \omega_j:\texttt{Type}} \quad \sigma:(\iota, \tuple{\eta_1:\omega_1,...,\eta_m:\omega_m}) \\*
		\Gamma\vdash \forall_{i \in [1,n]}{\tau_i \narr \omega_i}
	}{
		\Gamma\vdash \theta \narr \rho
	}
\]

A tuple type $\theta$ is narrower than the struct type $\sigma$ if the first $n$
fields of $\theta$ have a type $\tau$ narrower than the first $n$ fields type
$\omega$ of $\sigma$ (struct name and fields names are ignored)

\vspace{\baselineskip}

\subparagraph{Record}

\[
	\infrule{
		\Gamma\vdash n \geq 0 \quad \forall_{i\in[1,n]}{\eta_i \in \mathcal{N} \subseteq \mathcal{M}, \tau_i:\texttt{Type}} \quad \rho_1:\tuple{\eta_1:\tau_1,...,\eta_n:\tau_n} \\*
		\Gamma\vdash m \geq n \quad \forall_{j\in[1,m]}{\xi_j \in \mathcal{M} \subseteq \mathcal{I}, \omega_i:\texttt{Type}} \quad \rho_2:\tuple{\xi_1:\omega_1,...,\xi_m:\omega_m} \\*
		\Gamma\vdash \forall_{\eta':\tau'}{\exists_{\xi':\omega'}{\eta'=\xi' \land \tau' \narr \omega'}}
	}{
		\Gamma\vdash \rho_1 \narr \rho_2
	}
\]

A record type $\rho_1$ is narrower than a record type $\rho_2$ if, for all
fields $\eta'$ of type $\tau'$ in $\rho_1$, there is a field with the same name
$\xi'$ of type $\omega'$ in $\rho_2$, where $\tau'$ is narrower than $\omega'$

\vspace{\baselineskip}

\[
	\infrule{
		\Gamma\vdash n \geq 0 \quad \forall_{i\in[1,n]}{\eta \in \mathcal{N} \subseteq \mathcal{M}, \tau_i:\texttt{Type}} \quad \rho:\tuple{\eta_1:\tau_1,...,\eta_n:\tau_n} \\*
		\Gamma\vdash \epsilon \in \mathcal{I} \quad m \geq n \quad \forall_{j\in[1,m]}{\xi_j \in \mathcal{I}, \omega_j:\texttt{Type}} \quad \sigma:(\epsilon, \tuple{\xi_1:\omega_1,...,\xi_m:\omega_m}) \\*
		\Gamma\vdash \forall_{\eta':\tau'}{\exists_{\xi':\omega'}{\eta'=\xi' \land \tau' \narr \omega'}}
	}{
		\Gamma\vdash \rho \narr \phi
	}
\]

A record type $\rho$ is narrower than a struct type $\sigma$ if, for all fields
$\eta'$ of type $\tau'$ in $\rho$, there is a field with the same name $\xi'$
of type $\omega'$ in $\rho$, where $\tau'$ is narrower than $\omega'$ (struct
name is ignored)
\vspace{\baselineskip}

\subparagraph{Struct}

\[
	\infrule{
		\Gamma\vdash \iota \in \mathcal{I} \quad n \geq 0 \quad \forall_{i\in[1,n]}{\eta_i \in \mathcal{I}, \tau_i:\texttt{Type}} \quad \sigma_1:(\iota, \tuple{\eta_1:\tau_1,...,\eta_n:\tau_n}) \\*
		\Gamma\vdash \epsilon \in \mathcal{I} \quad m \geq n \quad \forall_{j\in[1,m]}{\xi_j \in \mathcal{I}, \omega_j:\texttt{Type}} \quad \sigma_2:(\epsilon, \tuple{\xi_1:\omega_1,...,\xi_m:\omega_m}) \\*
		\Gamma\vdash \iota = \epsilon
	}{
		\Gamma\vdash \sigma_1 \narr \sigma_2
	}
\]

A struct type $\sigma_1$ is narrower than a struct type $\sigma_2$ if and only if
they have the same name (fields are ignored since they're equal by
construction).

\comment{TODO: questa regola è ridondante?}
\vspace{\baselineskip}

\paragraph{Function}

\[
	\infrule{
		\Gamma\vdash \rho_1:\texttt{Type} \quad n \geq 0 \quad \forall_{i\in[1,n]}{\eta_i \in \mathcal{N} \subseteq \mathcal{M}, \tau_i:\texttt{Type}} \\*
		\phi_1:(\eta_1:\tau_1,...,\eta_n:\tau_n) \to \rho_1 \\*
		\Gamma\vdash \rho_2:\texttt{Type} \quad m \geq n  \quad \forall_{j\in[1,m]}{\xi_j \in \mathcal{M} \subseteq \mathcal{I}, \omega_i:\texttt{Type}} \\*
		\phi_2:(\xi_1:\omega_1,...,\xi_m:\omega_m) \to \rho_1 \\*
		\Gamma\vdash \rho_1 \narr \rho_2 \\*
		\Gamma\vdash \forall_{\eta':\tau'}{\exists_{\xi':\omega'}{\eta'=\xi' \land \tau' \narr \omega'}}
	}{
		\Gamma\vdash \phi_1 \narr \phi_2
	}
\]

A function type $\phi_1$ is narrower than a function type $\phi_2$ if the
return type $\rho_1$ is narrower than the other return type $\rho_2$, and if,
for all parameters $\eta'$ of type $\tau'$ in $\phi_1$, there is a parameter with
the same name $\xi'$ of type $\omega'$ in $\phi_2$, where $\tau'$ is narrower than
$\omega'$.

\comment{TODO: è gia un casino cosi :D, ma includere anche opzionalità dei parametri}

\comment{Eventualmente estrarre una regola per il singolo parametro}

\newpage

\section{Literals}

\subsection{Primitive types}

\subsubsection{Bool}

\[
	\Gamma\vdash \code{true}:\texttt{Bool} \qquad \Gamma \vdash \code{false}:\texttt{Bool}
\]

\subsubsection{Int}

\[
	\infrule{
		\Gamma\vdash x \in \mathbb{Z} \quad x \in [-2^{31}, 2^{31} - 1]
	}{
		\Gamma\vdash \code{$x$} : \texttt{Int} \\*
		\land \Gamma\vdash \code{$x$i} : \texttt{Int}
	}
\]

% \subparagraph{Examples} \verb|0, 0i, 69, -420, 32i, -99i|

\subsubsection{Natural}

\[
	\infrule{
		\Gamma\vdash x \in \mathbb{N} \quad x \in [0, 2^{31}]
	}{
		\Gamma\vdash \code{$x$n} : \texttt{Natural}
	}
\]

% \subparagraph{Examples} \verb|0n, 69n, 420n|

\subsubsection{Real}

\[
	\infrule{
		\Gamma\vdash x \in \mathbb{R}
	}{
		\Gamma\vdash \code{$x$} : \texttt{Real} \\*
		\land \Gamma\vdash \code{$x$r} : \texttt{Real}
	}
\]

% \subparagraph{Examples} \verb|0r, 0.0, 6.9, -4.20|

\subsubsection{Text}

\[
	\infrule{
		\Gamma\vdash x \in \text{Unicode} \setminus \set{"\backslash n", "\backslash r"} \quad n \geq 0
	}{
		\Gamma\vdash \code{"$x_1...x_n$"} : \texttt{Text} \\*
		\land \Gamma\vdash \code{'$x_1...x_n$'} : \texttt{Text}
	}
\]

\[
	\infrule{
		\Gamma\vdash x \in \text{Unicode} \quad n \geq 0
	}{
		\Gamma\vdash \code{"\{$x_1...x_n$\}"} : \texttt{Text} \\*
		\land \Gamma\vdash \code{'\{$x_1...x_n$\}'} : \texttt{Text}
	}
\]

% \subparagraph{Examples} \verb|"hello", 'world', "{hello}", '{world}'|

\comment{TODO: specificare il set $Unicode$ con i caratteri permessi da Java}

\subsection{Composite types}

\subsubsection{List}

\[
	\Gamma\vdash \code{[]}:\exists\texttt{Void}.\texttt{List}
\]

\[
	\infrule{
		\Gamma\vdash \tau:\texttt{Type} \qquad x:\tau \qquad n \in \mathbb{N}
	}{
		\Gamma\vdash \code{[$x_1, ...,x_n$]}:\exists\tau.\texttt{List}
	}
\]

\[
	\infrule{
		\Gamma\vdash n \in \mathbb{N} \quad \forall_{i \in [1,n]}{\tau_i:\texttt{Type}, x_i:\tau_i} \quad \omega = \bigsqcap_{i=1}^{n}{\tau_i}
	}{
		\Gamma\vdash \code{[$x_1, ...,x_n$]}:\exists\omega.\texttt{List}
	}
\]

List of $\omega$ (widest among $\tau$)

\subsubsection{Map}

\[
	\Gamma\vdash \code{[:]}:\exists\texttt{Void}.\exists\texttt{Void}.\texttt{Map}
\]

\[
	\infrule{
		\Gamma\vdash \kappa:\texttt{Type}
		\quad \nu:\texttt{Type}
		\quad \forall_{i\in[1,n]}{k_i:\kappa,v_i:\nu}
		\quad n \in \mathbb{N}
	}{
		\Gamma\vdash \code{[$k_1:v_1, ..., k_n:v_n$]}:\exists\kappa.\exists\nu.\texttt{Map}
	}
\]

\[
	\infrule{
		\Gamma\vdash \forall_{i \in [1,n]}{\kappa_i:\texttt{Type}, k_i:\kappa_i} \quad n \in \mathbb{N} \\*
		\Gamma\vdash \forall_{i \in [1,n]}{\nu_i:\texttt{Type}, v_i:\nu_i} \hfill\\*
		\Gamma\vdash \omega_\kappa = \bigsqcap_{i=1}^{n}{\kappa_i} \quad \omega_\nu = \bigsqcap_{i=1}^{n}{\nu_i}
	}{
		\Gamma\vdash \code{[$k_1:v_1, ..., k_n:v_n$]}:\exists\omega_\kappa.\exists\omega_\nu.\texttt{Map}
	}
\]

Map of $\omega_\kappa$ keys (widest  among $\kappa$) and $\omega_\nu$ (widest among $\nu$)

\subsection{Product types}

\subsubsection{Tuple}

\[
	\infrule{
		\Gamma\vdash n \geq 1 \quad \forall_{i\in[1,n]}{\tau_i:\texttt{Type}, x_i:\tau_i}
	}{
		\Gamma\vdash \tuplecode{$x_1,...,x_n$}:\tau_1 \times ... \times \tau_n
	}
\]

\subsubsection{Record}

\[
	\infrule{
		\Gamma\vdash n \geq 1 \quad \forall_{i\in[1,n]}{\tau_i:\texttt{Type}, x_i:\tau_i, a_i \in \mathcal{I}}
	}{
		\Gamma\vdash \tuplecode{$a_1:x_1,...,a_n:x_n$}:\tuple{a_1:\tau_1,...,a_2:\tau_n}
	}
\]

\subsubsection{Struct}

\comment{TODO: ...}

\subsection{Range}

\[
	\infrule{
		\Gamma\vdash \beta \in \texttt{Number} \quad b:\beta\\*
		\Gamma\vdash \epsilon \in \texttt{Number} \quad e:\epsilon \\*
		\Gamma\vdash \omega = \beta \sqcap \epsilon
	}{
		\Gamma\vdash \code{$b$..$e$} : \exists\omega.\texttt{Range}
	}
\]
Range from $b$ (begin) to $e$ (end)

\[
	\infrule{
		\Gamma\vdash \beta \in \texttt{Number} \quad b:\beta\\*
		\Gamma\vdash \epsilon \in \texttt{Number} \quad e:\epsilon \\*
		\Gamma\vdash \sigma \in \texttt{Number} \quad s:\sigma \\*
		\Gamma\vdash \omega = \beta \sqcap \epsilon \sqcap \sigma
	}{
		\Gamma\vdash \code{$b$..$e$..$s$} : \exists\omega.\texttt{Range}
	}
\]

Range from $b$ (begin) to $e$ (end) with a step $s$

\chapter{Operators}

\section{Operators}

\comment{TODO: aggiungere operatore dotref}

\begin{longtable}{cc}
	OPERATOR                        & SYMBOL     \\\hline
	PLUS                            & \verb$+$   \\\hline
	MINUS                           & \verb$-$   \\\hline
	MULTIPLICATION                  & \verb$*$   \\\hline
	DIVISION                        & \verb$/$   \\\hline
	MODULO                          & \verb$%$  \\\hline
	EXPONENTIAL                     & \verb$^$   \\\hline
	EQUAL                           & \verb$=$   \\\hline
	NOT EQUAL                       & \verb$!=$  \\\hline
	LESS EQUAL                      & \verb$<=$  \\\hline
	LESS                            & \verb$<$   \\\hline
	GREATER                         & \verb$>$   \\\hline
	GREATER EQUAL                   & \verb$>=$  \\\hline
	\hdashline\newpage\hdashline\hline
	CONTAINS                        & \verb$=~$  \\\hline
	IS CONTAINED                    & \verb$~=$  \\\hline
	\multirow{2}{*}{BOOL NOT}       & \verb$not$ \\
	                                & \verb$!$ \\\hline
	\multirow{2}{*}{BOOL AND}       & \verb$and$ \\
	                                & \verb$&$ \\\hline
	\multirow{2}{*}{BOOL OR}        & \verb$or$ \\
	                                & \verb$|$ \\\hline
	BITWISE NOT                     & \verb$!!$ \\\hline
	BITWISE AND                     & \verb$&&$ \\\hline
	BITWISE OR                      & \verb$||$ \\\hline
	BITWISE XOR                     & \verb$|||$ \\\hline
	BITWISE LEFT SHIFT              & \verb${<}$ \\\hline
	BITWISE RIGHT SHIFT             & \verb${>}$ \\\hline
	BITWISE RIGHT UNSIGNED SHIFT    & \verb${>>}$ \\\hline
	IS (type check)                 & \verb$is$  \\\hline
	AS (type coercion)              & \verb$as$  \\\hline
\end{longtable}

\section{Pseudo operators}

\begin{longtable}{rl}
	OPERATOR   & SYMBOL     \\\hline
	DOT ACCESS & \verb$.$ \\\hline
	AT ACCESS  & \verb$@$ \\\hline
\end{longtable}

\newpage

\section{Inference rules}

\[
	\infrule{
		\Gamma\vdash \tau:\texttt{Type} \quad l:\tau \\*
		\Gamma\vdash \sigma:\texttt{Type} \quad \tau \narr \sigma \quad r:\sigma
	}{
		\Gamma\vdash \code{$l$ = $r$} : \texttt{Bool}
	}
\]

\[
	\infrule{
		\Gamma\vdash \tau:\texttt{Type} \quad l:\tau \\*
		\Gamma\vdash \sigma:\texttt{Type} \quad r:\sigma \\*
		\Gamma\vdash \code{$l$ = $r$} : \texttt{Bool}
	}{
		\Gamma\vdash \code{$l$ != $r$} : \texttt{Bool}
	}
\]

If the "equal" operator is defined, the "not equal" operator is also defined

\vspace{\baselineskip}

\[
	\infrule{
		\Gamma\vdash \tau:\texttt{Type} \quad l:\tau \\*
		\Gamma\vdash \sigma:\texttt{Type} \quad r:\sigma \\*
		\Gamma\vdash \code{$l$ =\~{} $r$}: \texttt{Bool}
	}{
		\Gamma\vdash \code{$r$ \~{}= $l$}: \texttt{Bool}
	}
\]

If the "contains" operator is defined, the "contained" operator is also defined

\vspace{\baselineskip}

\[
	\texttt{Number} = \set{\texttt{Natural}, \texttt{Int}, \texttt{Real}}
\]

\[
	\infrule{
		\Gamma\vdash l:\texttt{Number} \quad \tau \in \texttt{Number}
	}{
		\Gamma\vdash \code{$l$ as $\tau$}: \tau
	}
\]

\comment{TODO(implementazione): -69 as Natural fallisce (out of bounds).
Dovrebbe?}

\[
	\infrule{
		\Gamma\vdash \tau \in \texttt{Number} \quad x:\tau
	}{
		\Gamma\vdash \code{+$x$}: \tau
	}
\]

\[
	\infrule{
		\Gamma\vdash \tau \in \texttt{Number}\setminus\set{\texttt{Natural}} \quad x:\tau
	}{
		\Gamma\vdash \code{-$x$}: \tau
	}
\]

\[
	\infrule{
		\Gamma\vdash \tau \in \texttt{Number} \quad l:\tau \quad r:\tau\\*
		op = \set{\code{=}, \code{!=}, \code{<}, \code{<=}, \code{>=}, \code{>=}, \code{>}}
	}{
		\Gamma\vdash \code{$l$ $op$ $r$}: \texttt{Bool}
	}
\]

\subsection{Primitive types}

\subsubsection{Void}

\comment{TODO: ?}

\subsubsection{Bool}

\[
	\infrule{
		\Gamma\vdash x:\texttt{Bool}
	}{
		\Gamma\vdash \code{!$x$}:\texttt{Bool}
	}
\]

\[
	\infrule{
		\Gamma\vdash l:\texttt{Bool} \quad r:\texttt{Bool} \\*
		op = \set{\code{and}, \code{\&}, \code{or}, \code{|}}
	}{
		\Gamma\vdash \code{$l$ $op$ $r$}:\texttt{Bool}
	}
\]

\subsubsection{Natural}

\paragraph{Plus}

\[
	\infrule{
		\Gamma\vdash l:\texttt{Natural} \quad r:\texttt{Natural}
	}{
		\Gamma\vdash \code{$l$ + $r$} : \texttt{Natural}
	}
\]

\[
	\infrule{
		\Gamma\vdash l:\texttt{Natural} \quad r:\texttt{Int}
	}{
		\Gamma\vdash \code{$l$ + $r$} : \texttt{Int}
	}
\]

\[
	\infrule{
		\Gamma\vdash l:\texttt{Natural} \quad r:\texttt{Real}
	}{
		\Gamma\vdash \code{$l$ + $r$} : \texttt{Real}
	}
\]

\paragraph{Minus}

\[
	\infrule{
		\Gamma\vdash l:\texttt{Natural} \quad r:\texttt{Natural}
	}{
		\Gamma\vdash \code{$l$ - $r$} : \texttt{Int}
	}
\]

\comment{TODO: specificare se $l-r\geq0$}

\[
	\infrule{
		\Gamma\vdash l:\texttt{Natural} \quad r:\texttt{Int}
	}{
		\Gamma\vdash \code{$l$ - $r$} : \texttt{Int}
	}
\]

\[
	\infrule{
		\Gamma\vdash l:\texttt{Natural} \quad r:\texttt{Real}
	}{
		\Gamma\vdash \code{$l$ - $r$} : \texttt{Real}
	}
\]

\paragraph{Multiplication}

\[
	\infrule{
		\Gamma\vdash l:\texttt{Natural} \quad r:\texttt{Natural}
	}{
		\Gamma\vdash \code{$l$ * $r$} : \texttt{Natural}
	}
\]


\[
	\infrule{
		\Gamma\vdash l:\texttt{Natural} \quad r:\texttt{Int}
	}{
		\Gamma\vdash \code{$l$ * $r$} : \texttt{Int}
	}
\]

\[
	\infrule{
		\Gamma\vdash l:\texttt{Natural} \quad r:\texttt{Real}
	}{
		\Gamma\vdash \code{$l$ * $r$} : \texttt{Real}
	}
\]

\paragraph{Division}

\[
	\infrule{
		\Gamma\vdash l:\texttt{Natural} \quad r:\texttt{Natural}
	}{
		\Gamma\vdash \code{$l$ / $r$} : \texttt{Real}
	}
\]

\[
	\infrule{
		\Gamma\vdash l:\texttt{Natural} \quad r:\texttt{Int}
	}{
		\Gamma\vdash \code{$l$ / $r$} : \texttt{Real}
	}
\]

\[
	\infrule{
		\Gamma\vdash l:\texttt{Natural} \quad r:\texttt{Real}
	}{
		\Gamma\vdash \code{$l$ / $r$} : \texttt{Real}
	}
\]

\paragraph{Modulo}

\[
	\infrule{
		\Gamma\vdash l:\texttt{Natural} \quad r:\texttt{Natural}
	}{
		\Gamma\vdash \code{$l$ \% $r$} : \texttt{Natural}
	}
\]

\[
	\infrule{
		\Gamma\vdash l:\texttt{Natural} \quad r:\texttt{Int}
	}{
		\Gamma\vdash \code{$l$ \% $r$} : \texttt{Int}
	}
\]

\[
	\infrule{
		\Gamma\vdash l:\texttt{Natural} \quad r:\texttt{Real}
	}{
		\Gamma\vdash \code{$l$ \% $r$} : \texttt{Real}
	}
\]

\paragraph{Power}

\[
	\infrule{
		\Gamma\vdash l:\texttt{Natural} \quad r:\texttt{Natural}
	}{
		\Gamma\vdash \code{$l$ \^{} $r$} : \texttt{Natural}
	}
\]

\[
	\infrule{
		\Gamma\vdash l:\texttt{Natural} \quad r:\texttt{Int}
	}{
		\Gamma\vdash \code{$l$ \^{} $r$} : \texttt{Int}
	}
\]

\[
	\infrule{
		\Gamma\vdash l:\texttt{Natural} \quad r:\texttt{Real}
	}{
		\Gamma\vdash \code{$l$ \^{} $r$} : \texttt{Real}
	}
\]

\subsubsection{Int}

\paragraph{Bitwise not}

\[
	\infrule{
		\Gamma\vdash x:\texttt{Int}
	}{
		\Gamma\vdash \code{!!$x$} : \texttt{Int}
	}
\]

\paragraph{Bitwise and}

\[
	\infrule{
		\Gamma\vdash l:\texttt{Int} \quad r:\texttt{Natural}
	}{
		\Gamma\vdash \code{$l$ \&\& $r$} : \texttt{Int}
	}
\]

\[
	\infrule{
		\Gamma\vdash l:\texttt{Int} \quad r:\texttt{Int}
	}{
		\Gamma\vdash \code{$l$ \&\& $r$} : \texttt{Int}
	}
\]

\paragraph{Bitwise or}

\[
	\infrule{
		\Gamma\vdash l:\texttt{Int} \quad r:\texttt{Natural}
	}{
		\Gamma\vdash \code{$l$ || $r$} : \texttt{Int}
	}
\]

\[
	\infrule{
		\Gamma\vdash l:\texttt{Int} \quad r:\texttt{Int}
	}{
		\Gamma\vdash \code{$l$ || $r$} : \texttt{Int}
	}
\]

\paragraph{Bitwise xor}

\[
	\infrule{
		\Gamma\vdash l:\texttt{Int} \quad r:\texttt{Natural}
	}{
		\Gamma\vdash \code{$l$ ||| $r$} : \texttt{Int}
	}
\]

\[
	\infrule{
		\Gamma\vdash l:\texttt{Int} \quad r:\texttt{Int}
	}{
		\Gamma\vdash \code{$l$ ||| $r$} : \texttt{Int}
	}
\]

\paragraph{Bitwise left shift}

\[
	\infrule{
		\Gamma\vdash l:\texttt{Int} \quad r:\texttt{Natural}
	}{
		\Gamma\vdash \code{$l$ \{<\} $r$} : \texttt{Int}
	}
\]

\[
	\infrule{
		\Gamma\vdash l:\texttt{Int} \quad r:\texttt{Int}
	}{
		\Gamma\vdash \code{$l$ \{<\} $r$} : \texttt{Int}
	}
\]

\paragraph{Bitwise right shift}

\[
	\infrule{
		\Gamma\vdash l:\texttt{Int} \quad r:\texttt{Natural}
	}{
		\Gamma\vdash \code{$l$ \{>\} $r$} : \texttt{Int}
	}
\]

\[
	\infrule{
		\Gamma\vdash l:\texttt{Int} \quad r:\texttt{Int}
	}{
		\Gamma\vdash \code{$l$ \{>\} $r$} : \texttt{Int}
	}
\]

\paragraph{Bitwise right unsigned shift}

\[
	\infrule{
		\Gamma\vdash l:\texttt{Int} \quad r:\texttt{Natural}
	}{
		\Gamma\vdash \code{$l$ \{>>\} $r$} : \texttt{Int}
	}
\]

\[
	\infrule{
		\Gamma\vdash l:\texttt{Int} \quad r:\texttt{Int}
	}{
		\Gamma\vdash \code{$l$ \{>>\} $r$} : \texttt{Int}
	}
\]

\paragraph{Plus}

\[
	\infrule{
		\Gamma\vdash l:\texttt{Int} \quad r:\texttt{Natural}
	}{
		\Gamma\vdash \code{$l$ + $r$} : \texttt{Int}
	}
\]

\[
	\infrule{
		\Gamma\vdash l:\texttt{Int} \quad r:\texttt{Int}
	}{
		\Gamma\vdash \code{$l$ + $r$} : \texttt{Int}
	}
\]

\[
	\infrule{
		\Gamma\vdash l:\texttt{Int} \quad r:\texttt{Real}
	}{
		\Gamma\vdash \code{$l$ + $r$} : \texttt{Real}
	}
\]

\paragraph{Minus}

\[
	\infrule{
		\Gamma\vdash l:\texttt{Int} \quad r:\texttt{Natural}
	}{
		\Gamma\vdash \code{$l$ - $r$} : \texttt{Int}
	}
\]

\[
	\infrule{
		\Gamma\vdash l:\texttt{Int} \quad r:\texttt{Int}
	}{
		\Gamma\vdash \code{$l$ - $r$} : \texttt{Int}
	}
\]

\[
	\infrule{
		\Gamma\vdash l:\texttt{Int} \quad r:\texttt{Real}
	}{
		\Gamma\vdash \code{$l$ - $r$} : \texttt{Real}
	}
\]

\paragraph{Multiplication}

\[
	\infrule{
		\Gamma\vdash l:\texttt{Int} \quad r:\texttt{Natural}
	}{
		\Gamma\vdash \code{$l$ * $r$} : \texttt{Int}
	}
\]

\[
	\infrule{
		\Gamma\vdash l:\texttt{Int} \quad r:\texttt{Int}
	}{
		\Gamma\vdash \code{$l$ * $r$} : \texttt{Int}
	}
\]

\[
	\infrule{
		\Gamma\vdash l:\texttt{Int} \quad r:\texttt{Real}
	}{
		\Gamma\vdash \code{$l$ * $r$} : \texttt{Real}
	}
\]

\paragraph{Division}

\[
	\infrule{
		\Gamma\vdash l:\texttt{Int} \quad r:\texttt{Natural}
	}{
		\Gamma\vdash \code{$l$ / $r$} : \texttt{Real}
	}
\]

\[
	\infrule{
		\Gamma\vdash l:\texttt{Int} \quad r:\texttt{Int}
	}{
		\Gamma\vdash \code{$l$ / $r$} : \texttt{Real}
	}
\]

\[
	\infrule{
		\Gamma\vdash l:\texttt{Int} \quad r:\texttt{Real}
	}{
		\Gamma\vdash \code{$l$ / $r$} : \texttt{Real}
	}
\]

\paragraph{Modulo}

\[
	\infrule{
		\Gamma\vdash l:\texttt{Int} \quad r:\texttt{Natural}
	}{
		\Gamma\vdash \code{$l$ \% $r$} : \texttt{Int}
	}
\]

\[
	\infrule{
		\Gamma\vdash l:\texttt{Int} \quad r:\texttt{Int}
	}{
		\Gamma\vdash \code{$l$ \% $r$} : \texttt{Int}
	}
\]

\[
	\infrule{
		\Gamma\vdash l:\texttt{Int} \quad r:\texttt{Real}
	}{
		\Gamma\vdash \code{$l$ \% $r$} : \texttt{Real}
	}
\]

\paragraph{Power}

\[
	\infrule{
		\Gamma\vdash l:\texttt{Int} \quad r:\texttt{Natural}
	}{
		\Gamma\vdash \code{$l$ \^{} $r$} : \texttt{Real}
	}
\]

\[
	\infrule{
		\Gamma\vdash l:\texttt{Int} \quad r:\texttt{Int}
	}{
		\Gamma\vdash \code{$l$ \^{} $r$} : \texttt{Real}
	}
\]

\[
	\infrule{
		\Gamma\vdash l:\texttt{Int} \quad r:\texttt{Real}
	}{
		\Gamma\vdash \code{$l$ \^{} $r$} : \texttt{Real}
	}
\]

\subsubsection{Real}

\paragraph{Plus}

\[
	\infrule{
		\Gamma\vdash l:\texttt{Int} \quad \tau \in \texttt{Number} \quad r:\tau
	}{
		\Gamma\vdash \code{$l$ + $r$} : \texttt{Real}
	}
\]

\paragraph{Minus}

\[
	\infrule{
		\Gamma\vdash l:\texttt{Int} \quad \tau \in \texttt{Number} \quad r:\tau
	}{
		\Gamma\vdash \code{$l$ - $r$} : \texttt{Real}
	}
\]

\paragraph{Multiplication}

\[
	\infrule{
		\Gamma\vdash l:\texttt{Int} \quad \tau \in \texttt{Number} \quad r:\tau
	}{
		\Gamma\vdash \code{$l$ * $r$} : \texttt{Real}
	}
\]

\paragraph{Division}

\[
	\infrule{
		\Gamma\vdash l:\texttt{Int} \quad \tau \in \texttt{Number} \quad r:\tau
	}{
		\Gamma\vdash \code{$l$ / $r$} : \texttt{Real}
	}
\]

\paragraph{Modulo}

\[
	\infrule{
		\Gamma\vdash l:\texttt{Int} \quad \tau \in \texttt{Number} \quad r:\tau
	}{
		\Gamma\vdash \code{$l$ \% $r$} : \texttt{Real}
	}
\]

\paragraph{Power}

\[
	\infrule{
		\Gamma\vdash l:\texttt{Int} \quad \tau \in \texttt{Number} \quad r:\tau
	}{
		\Gamma\vdash \code{$l$ \^{} $r$} : \texttt{Real}
	}
\]

\subsubsection{Text}

\paragraph{As}

\[
	\infrule{
		\Gamma\vdash l:\texttt{Text} \quad \tau \in \texttt{Number} \cup \set{\texttt{Text}, \texttt{Bool}}
	}{
		\Gamma\vdash \code{$l$ as $\tau$}: \tau
	}
\]

\paragraph{Plus}

\[
	\infrule{
		\Gamma\vdash l:\texttt{Text} \quad \tau:\texttt{Type} \quad r:\tau
	}{
		\Gamma\vdash \code{$l$ + $r$}: \texttt{Text}
	}
\]

\paragraph{Minus}

\[
	\infrule{
		\Gamma\vdash l:\texttt{Text} \quad r:\texttt{Text}
	}{
		\Gamma\vdash \code{$l$ - $r$}: \texttt{Text}
	}
\]

\paragraph{Power}

\[
	\infrule{
		\Gamma\vdash l:\texttt{Text} \quad \tau \in \texttt{Number} \quad r:\tau
	}{
		\Gamma\vdash \code{$l$ \^{} $r$}: \texttt{Text}
	}
\]

\comment{TODO(implementazione): non dovremmo permettere Real}

\paragraph{Division}

\[
	\infrule{
		\Gamma\vdash l:\texttt{Text} \quad r:\texttt{Text}
	}{
		\Gamma\vdash \code{$l$ / $r$}: \exists\texttt{Text}.\texttt{List}
	}
\]

\paragraph{Equal}

\[
	\infrule{
		\Gamma\vdash l:\texttt{Text} \quad r:\texttt{Text}
	}{
		\Gamma\vdash \code{$l$ = $r$}: \texttt{Bool}
	}
\]

\paragraph{Less}

\[
	\infrule{
		\Gamma\vdash l:\texttt{Text} \quad r:\texttt{Text}
	}{
		\Gamma\vdash \code{$l$ < $r$}: \texttt{Bool}
	}
\]

\paragraph{Less equal}

\[
	\infrule{
		\Gamma\vdash l:\texttt{Text} \quad r:\texttt{Text}
	}{
		\Gamma\vdash \code{$l$ <= $r$}: \texttt{Bool}
	}
\]

\paragraph{Greater equal}

\[
	\infrule{
		\Gamma\vdash l:\texttt{Text} \quad r:\texttt{Text}
	}{
		\Gamma\vdash \code{$l$ >= $r$}: \texttt{Bool}
	}
\]

\paragraph{Greater}

\[
	\infrule{
		\Gamma\vdash l:\texttt{Text} \quad r:\texttt{Text}
	}{
		\Gamma\vdash \code{$l$ > $r$}: \texttt{Bool}
	}
\]

\paragraph{Contains}

\[
	\infrule{
		\Gamma\vdash l:\texttt{Text} \quad r:\texttt{Text}
	}{
		\Gamma\vdash \code{$l$ =\~{} $r$}: \texttt{Bool}
	}
\]

\paragraph{Access}

\[
	\text{Index access}\quad
	\infrule{
		\Gamma\vdash l:\texttt{Text} \quad r:\texttt{Int}
	}{
		\Gamma\vdash \code{$l$ @ $r$}: \texttt{Text}
	}
\]

\[
	\text{"Substring"}\quad
	\infrule{
		\Gamma\vdash l:\texttt{Text} \quad r:\exists\texttt{Int}.\texttt{Range}
	}{
		\Gamma\vdash \code{$l$ @ $r$}: \texttt{Text}
	}
\]

\comment{TODO: specificare parametri range}

\[
	\text{Count access}\quad
	\infrule{
		\Gamma\vdash l:\texttt{Text}
	}{
		\Gamma\vdash \code{$l$ @ \%}: \texttt{Int}
	}
\]

\[
	\text{Range access}\quad
	\infrule{
		\Gamma\vdash l:\texttt{Text}
	}{
		\Gamma\vdash \code{$l$ @ @}: \exists\texttt{Int}.\texttt{Range}
	}
\]

\[
	\text{Values access}\quad
	\infrule{
		\Gamma\vdash l:\texttt{Text}
	}{
		\Gamma\vdash \code{$l$ @ *}: \exists\texttt{Text}.\texttt{List}
	}
\]

\comment{TODO: legacy access syntax. Rimuovere?}

\subsection{Composite types}

\subsubsection{List}

\paragraph{As}

\[
	\infrule{
		\Gamma\vdash \tau:\texttt{Type} \quad l:\exists\tau.\texttt{List} \\*
		\Gamma\vdash \rho:\texttt{List}[\sigma:=\texttt{Type}] \\*
		\Gamma\vdash \code{x:$\tau$ as $\sigma$} : \sigma
	}{
		\Gamma\vdash \code{$l$ as $\rho$} : \rho
	}
\]

A term of type $\tau$ can be casted to type $\sigma$ $\implies$ A list of
$\tau$ can be casted to the type $\rho$

\paragraph{Contains}

\[
	\infrule{
		\Gamma\vdash \tau:\texttt{Type} \quad l:\exists\tau.\texttt{List} \\*
		\Gamma\vdash \sigma:\texttt{Type} \quad \sigma \narr \tau \quad r:\sigma
	}{
		\Gamma\vdash \code{$l$ =\~{} $r$} : \texttt{Bool}
	}
\]

A list of $\tau$ contains an element of type $\sigma$ (narrower than $\tau$)

\[
	\infrule{
		\Gamma\vdash \tau:\texttt{Type} \quad l:\exists\tau.\texttt{List}\\*
		\Gamma\vdash \sigma:\texttt{Type} \quad \sigma \narr \tau \quad r:\exists\sigma.\texttt{List}
	}{
		\Gamma\vdash \code{$l$ =\~{} $r$} : \texttt{Bool}
	}
\]

A list of $\tau$ contains the elements of a list of $\sigma$ (narrower than
$\tau$)

\paragraph{Access}

\[
	\text{Index access}\quad
	\infrule{
		\Gamma\vdash \tau:\texttt{Type}
		\quad l:\exists\tau.\texttt{List}
		\quad r:\texttt{Int}
	}{
		\Gamma\vdash \code{$l$ @ $r$}: \tau
	}
\]

\[
	\infrule{
		\Gamma\vdash \tau:\texttt{List}[\sigma:=\texttt{Type}]
		\quad l:\tau
		\quad r:\exists\texttt{Int}.\texttt{Range}
	}{
		\Gamma\vdash \code{$l$ @ $r$}: \tau
	}
\]

\[
	\text{Count access}\quad
	\infrule{
		\Gamma\vdash \tau:\texttt{Type} \quad l:\exists\tau.\texttt{List}
	}{
		\Gamma\vdash \code{$l$ @ \%}: \texttt{Int}
	}
\]

\[
	\text{Range access}\quad
	\infrule{
		\Gamma\vdash l:\texttt{Text}
	}{
		\Gamma\vdash \code{$l$ @ @}: \exists\texttt{Int}.\texttt{Range}
	}
\]

\comment{TODO: operatori aritmetici (+ concat, - remove) e uguaglianza}

\subsubsection{Map}

\paragraph{As}

\[
	\infrule{
		\Gamma\vdash \tau:\texttt{Map}[\kappa_1:=\texttt{Type},\nu_1:=\texttt{Type}] \quad l:\tau \\*
		\Gamma\vdash \rho:\texttt{Map}[\kappa_2:=\texttt{Type},\nu_2:=\texttt{Type}] \\*
		\Gamma\vdash \code{k:$\kappa_1$ as $\kappa_2$} : \kappa_2 \land \code{v:$\nu_1$ as $\nu_2$} : \nu_2
	}{
		\Gamma\vdash \code{$l$ as $\rho$} : \rho
	}
\]

A term of type $\kappa_1$ and be casted to type $\kappa_2$ $\land$ A term of
type $\nu_1$ can be casted to type $\nu_2$ $\implies$ A map of type $\tau$ can
be casted to the type $\rho$


\paragraph{Contains}

\[
	\infrule{
		\Gamma\vdash \tau:\texttt{Map}[\kappa:=\texttt{Type},\nu:=\texttt{Type}] \quad l:\tau \\*
		\Gamma\vdash \sigma:\texttt{Type} \quad \sigma \narr \kappa \quad r:\sigma
	}{
		\Gamma\vdash \code{$l$ =\~{} $r$} : \texttt{Bool}
	}
\]

A map $l$ of type $\tau$ contains a key of type $\sigma$ (narrower than
$\kappa$)

\[
	\infrule{
		\Gamma\vdash \tau:\texttt{Map}[\kappa_1:=\texttt{Type},\nu_1:=\texttt{Type}] \quad l:\tau \\*
		\Gamma\vdash \rho:\texttt{Map}[\kappa_2:=\texttt{Type},\nu_2:=\texttt{Type}] \quad r:\rho \\*
		\Gamma\vdash \kappa_2 \narr \kappa_1 \land \nu_2 \narr \nu_1
	}{
		\Gamma\vdash \code{$l$ =\~{} $r$} : \texttt{Bool}
	}
\]

A map $l$ of type $\tau$ contais the key-value pairs of type $\kappa_2$
(narrower than $\kappa_1$) and $\nu_2$ (narrower than $\nu_1$) of the map $r$
of type $\rho$

\paragraph{Access}

\[
	\text{Access by key}\quad
	\infrule{
		\Gamma\vdash \tau:\texttt{Map}[\kappa:=\texttt{Type},\nu:=\texttt{Type}] \quad l:\tau \\*
		\Gamma\vdash \rho:\texttt{Type} \quad \rho \narr \kappa \quad r:\rho
	}{
		\Gamma\vdash \code{$l$ @ $r$}: \nu
	}
\]

\[
	\text{Count access}\quad
	\infrule{
		\Gamma\vdash \tau:\texttt{Map}[\kappa:=\texttt{Type},\nu:=\texttt{Type}] \quad l:\tau \\*
	}{
		\Gamma\vdash \code{$l$ @ \%}: \texttt{Int}
	}
\]

\[
	\text{Keys access}\quad
	\infrule{
		\Gamma\vdash \tau:\texttt{Map}[\kappa:=\texttt{Type},\nu:=\texttt{Type}] \quad l:\tau \\*
	}{
		\Gamma\vdash \code{$l$ @ @}: \exists\kappa.\texttt{List}
	}
\]

\[
	\text{Values access}\quad
	\infrule{
		\Gamma\vdash \tau:\texttt{Map}[\kappa:=\texttt{Type},\nu:=\texttt{Type}] \quad l:\tau \\*
	}{
		\Gamma\vdash \code{$l$ @ @}: \exists\nu.\texttt{List}
	}
\]

\subsection{Product types}

\[
	\texttt{Products} = \set{\texttt{Tuple, Record, Struct}}
\]

\[
	\infrule{
		\Gamma\vdash \tau \in \texttt{Products} \quad l:\tau \\*
		\Gamma\vdash \sigma \in \texttt{Products} \quad r:\sigma
	}{
		\Gamma\vdash \code{$l$ = $r$}: \texttt{Bool}
	}
\]


\subsubsection{Tuple}

\paragraph{Access}

\[
	\text{"Slice" access}\quad
	\infrule{
		\Gamma\vdash n \in \mathbb{N} \quad \forall_{i\in[1,n]}{\tau_i:\texttt{Type}} \quad l:\tau_1 \times ... \times \tau_n \\*
		\Gamma\vdash r:\exists\texttt{Int}.\texttt{Range}
	}{
		\Gamma\vdash \code{$l$ @ $r$}: \bigtimes_{i \in r} \tau_i
	}
\]

Example: \code{<0,1,2,3,4,5> @ 2..4 = <2 3 4>}

With step: \code{<0,1,2,3,4,5> @ 0..5..2 = <0 2 4>}

Reverse: \code{<0,1,2,3,4,5> @ 4..2 = <4 3 2>}

\[
	\text{Access by index}\quad
	\infrule{
		\Gamma\vdash n \in \mathbb{N} \quad \forall_{i\in[1,n]}{\tau_i:\texttt{Type}} \quad l:\tau_1 \times ... \times \tau_n \\*
		\Gamma\vdash r:\texttt{Int} \quad r \in [1,n]
	}{
		\Gamma\vdash \code{$l$ @ $r$}: \tau_r
	}
\]

\[
	\infrule{
		\Gamma\vdash n \in \mathbb{N} \quad \forall_{i\in[1,n]}{\tau_i:\texttt{Type}} \quad l:\tau_1 \times ... \times \tau_n \\*
	}{
		\begin{aligned}
			\text{(Count access)} \qquad & \Gamma\vdash \code{$l$ @ \%}: \texttt{Int}                           \\
			\text{(Range access)} \qquad & \land\Gamma\vdash \code{$l$ @ @}: \exists\texttt{Int}.\texttt{Range} 
		\end{aligned}
	}
\]

\subsubsection{Record}

\paragraph{Access}

\[
	\text{Access by index}\quad
	\infrule{
		\Gamma\vdash n \geq 1 \quad \forall_{i\in[1,n]}{\tau_i:\texttt{Type}, x_i:\tau_i, a_i \in \mathcal{I}} \\*
		\Gamma\vdash l:\tuple{a_1:\tau_1,...,a_n:\tau_n}\\*
		\Gamma\vdash r:\texttt{Int} \quad r \in [1,n]
	}{
		\Gamma\vdash \code{$l$ @ $r$}: \tau_r
	}
\]

\comment{TODO(implemenatzione): permettere solo a Tuple in modo da rendere esplicito la perdita degli identificatori}

\[
	\infrule{
		\Gamma\vdash n \geq 1 \quad \forall_{i\in[1,n]}{\tau_i:\texttt{Type}, x_i:\tau_i, a_i \in \mathcal{I}} \\*
		\Gamma\vdash l:\tuple{a_1:\tau_1,...,a_n:\tau_n}\\*
	}{
		\begin{aligned}
			\text{(Count access)} \qquad & \Gamma\vdash \code{$l$ @ \%}: \texttt{Int}                     \\*
			\text{(Range access)} \qquad & \land\Gamma\vdash \code{$l$ @ @}: \exists\texttt{Int}.\texttt{Range}\\*
			\text{(Fields access)} \qquad & \land\Gamma\vdash \code{$l$ @ *}: \exists\texttt{Text}.\texttt{List}
		\end{aligned}
	}
\]

\paragraph{Dot access}

\[
	\text{Access by name}\quad
	\infrule{
		\Gamma\vdash n \geq 1 \quad \forall_{i\in[1,n]}{\tau_i:\texttt{Type}, x_i:\tau_i, a_i \in \mathcal{I}} \\*
		\Gamma\vdash l:\tuple{a_1:\tau_1,...,a_n:\tau_n}\\*
		\Gamma\vdash r \in \mathcal{I} \quad \exists_j a_j = r
	}{
		\Gamma\vdash \code{$l$ @ $r$}: \tau_j
	}
\]


\subsubsection{Struct}

\comment{TODO: ...}

\subsubsection{Function}

\comment{TODO: ...}

\chapter{Statements}

\section{Variable declaration}

\[
	\infrule{
		\Gamma\vdash \tau:\texttt{Type}
	}{
		\Gamma, \code{$a$ $\tau$} \vdash a:\tau
	}
\]

\comment{TODO: specificare il tipo Var come da implementazione?}

\section{Constant declaration}

\[
	\infrule{
		\Gamma\vdash \tau:\texttt{Type} \quad e:\tau
	}{
		\Gamma, \code{const $A$ $\tau$ <- $e$} \vdash A:\tau \qquad \text{(esplicit type)} \\*
		\land \Gamma, \code{const $A$ <- $e$} \vdash A:\tau \qquad \text{(implicit type)}
	}
\]

\section{For loop}

\[
	\infrule{
		\Gamma\vdash \tau:\texttt{Type} \quad l:\exists\tau.\texttt{List}
	}{
		\Gamma, \code{for $i$ in $l$} \vdash i:\tau
	}
\]

\[
	\infrule{
		\Gamma\vdash \kappa:\texttt{Type} \quad \nu:\texttt{Type} \quad m:\exists\kappa.\exists\nu.\texttt{Map}
	}{
		\Gamma, \code{for $k$:$v$ in $m$} \vdash k:\kappa \quad v:\nu
	}
\]

\comment{TODO: deconstruction}

\chapter{Operational semantics}

This section provides a semi-formal definition of the operational semantics of
the language, definining the state and context and how they mutate with every
expression.

\section{Notation and terminology}

Two environments will be used:

\begin{itemize}
	\item $\envLoc :=$ \textbf{location environment}. Describes the mappings
		between identifiers and their "location/address" in memory.
		\[
			\begin{split}
				\envLoc: \mathcal{I} & \to \mathbb{Z}\\*
				\envLoc(i) & \mapsto a
			\end{split}
		\]
		NOTE: "location" and "memory" are used loosely here. They do not
		represent an actual memory address, but rather an abstraction of it.
	\item $\envStore :=$ \textbf{store environment}. Describes the mappings
		between a location in memory and its data content
		\[
				\mathcal{V} = \set{v \mid v := \text{any value or data}}
		\]
		\[
			\begin{split}
				\envStore: \mathbb{Z} & \to \mathcal{V} \\*
				\envStore(a) & \mapsto v
			\end{split}
		\]
\end{itemize}

\paragraph{Environment extension} To denote a environment extending, modifying
or overriding an existing one, this notation will be used $\envLoc\env{i \mapsto
a}$ where it's defined as follows:

\[
	\envLoc\env{i \mapsto a}(j) \mapsto \begin{cases}
		a          &\text{ if } i = j\\
		\envLoc(j) &\text{ if } i \neq j\\
	\end{cases}
\]

Multiple mappings may be specified: $\envLoc\env{i_1 \mapsto a_1, i_2 \mapsto a_2}$

\paragraph{Environment modification} If the evaluation of an expression
results in an alteration of an environment, the judgment may be followed by a
"double turnstile" symbol $\vDash$ and the new environments.\\*
For example: $\envLoc,\envStore \vdash e \vDash \envStore_2$

\paragraph{Expression value} To denote to what value an expression evaluates
to, the $\equiv$ symbol will be used.\\*
For example: $\envLoc,\envStore \vdash e \equiv v \vDash \envStore_2$.\\*
To specify the type of the value, the notation $\tau(v)$ will be used, where
$\tau$ can be any \texttt{Type} and $v$ any value

\section{Operational rules}

\subsection{Literals}

\comment{TODO: come nelle regole di inferenza?}

\paragraph{Bool}

\[
	\oprule{
		b \in \set{\texttt{true}, \texttt{false}}
	}{
		\envLoc,\envStore\vdash \code{$b$} \equiv \texttt{Bool($b$)}
	}
\]

\paragraph{Natural}

\[
	\oprule{
		n := \text{ natural literal }
	}{
		\envLoc,\envStore\vdash \code{$n$} \equiv \texttt{Natural($n$)}
	}
\]

\paragraph{Int}

\[
	\oprule{
		i := \text{ integer literal }
	}{
		\envLoc,\envStore\vdash \code{$i$} \equiv \texttt{Int($i$)}
	}
\]

\paragraph{Real}

\[
	\oprule{
		r := \text{ real literal }
	}{
		\envLoc,\envStore\vdash \code{$r$} \equiv \texttt{Real($r$)}
	}
\]


\paragraph{Text}

\[
	\oprule{
		t := \text{ text literal }
	}{
		\envLoc,\envStore\vdash \code{$t$} \equiv \texttt{Text($t$)}
	}
\]

\comment{TODO: variables interpolation?}

\paragraph{Identifiers}

\[
	\oprule{
		i \in \mathcal{I}\\*
		\envLoc(i) = l\\*
		\envStore(l) = v
	}{
		\envLoc,\envStore\vdash \code{$i$} \equiv v
	}
\]

\[
	\oprule{
		\envLoc,\envStore\vdash e \equiv v_e \vDash \envLoc_2,\envStore_2 \\*
		i \in \mathcal{I} \\*
		\envLoc(i) = l \\*
		\envStore(l) = v_l
	}{
		\envLoc,\envStore\vdash \code{$i$ <- $e$} \vDash \envStore_2[l \mapsto v_e]
	}
\]

NOTE: evaluating $e$ may cause side effects (e.g. a function
invokation), hence the new enviroments $\envLoc_2$ and
$\envStore_2$



\end{document}


